\par\bigskip\Needspace{7\baselineskip}
\pdfbookmark[1]{Theorie der biquadratischen Reste. III}{nachlass-biquadratic-III}
\begin{center}[III.]\par\medskip\textsc{Theorie der biquadratischen Reste.}\end{center}
\BookRule{.08\textwidth}
\renewcommand{\Article}[1]{\par\medskip\Needspace{4\baselineskip}\pdfbookmark[2]{Art. #1}{nachlass-biquadratic-III-art-#1}\centerline{\textbf{#1.}}\nobreak\smallskip}
\Article{1}
Die kleinsten Reste zum Modul $m=a+bi$ sind die ganzen Zahlen $\mu=\alpha+\beta i$, für die bei $\frac{\mu}{m}=x+yi$ die Größen $x$ und $y$ positiv und kleiner als $1$ sind. Hinzu kommt der Rest $0$\begingroup\renewcommand{\thefootnote}{*)}\footnote{und, falls $a$ und $b$ einen gemeinsamen Teiler $e$ haben, die Zahlen $\frac{m}{e},\frac{2m}{e},\frac{3m}{e}\dotstwo\frac{(e-1)m}{e}$. Diesen Fall wollen wir jedoch vorerst aus der Untersuchung ausschließen.}\endgroup. Ihre Anzahl ist $=aa+bb$.

\Article{2}
Ist $aa+bb$ ungerade, so hat $aa+bb$ die Form $4n+1$. Lässt man den kleinsten Rest $0$ weg, so zerfallen die übrigen in vier Klassen. Zur ersten Klasse $f$ zählen wir diejenigen, bei denen $x$ und $y$ kleiner als $\frac12$ sind,
\SourceLeaf{342}
\[
\begin{array}{lcll}
\text{die zweite}&f'&\text{bei denen }x>\tfrac12,&y<\tfrac12\\
\text{dritte}&f''&x>\tfrac12,&y>\tfrac12\\
\text{vierte}&f'''&x<\tfrac12,&y>\tfrac12
\end{array}
\]
Man erhält alle Reste
\[
\begin{array}{rcl}
f'&\text{aus}&if+m\\
f''&\text{aus}&-f+(1+i)m\\
f'''&\text{aus}&-if+im
\end{array}
\]

\Needspace{27\baselineskip}\Article{3}
Es sei $M$ eine andere Zahl, die mit $m$ keinen Faktor gemeinsam hat. Dann gilt
\[
M^{aa+bb-1}\equiv1\ (\mathrm{mod.}\ m)
\]
folglich ist $M^{\frac14(aa+bb-1)}$ entweder $\equiv1$, oder $\equiv i$, oder $\equiv-1$, oder $\equiv-i$, das heißt $\equiv i^{\varepsilon}$, wobei $\varepsilon$ eine der vier Zahlen $0,1,2,3$ bezeichnet. Im ersten Fall ist $M$ ein biquadratischer Rest von $m$, also auch ein quadratischer. Im dritten Fall ist $M$ ein quadratischer, aber kein biquadratischer Rest; im zweiten und vierten Fall sowohl ein quadratischer als auch ein biquadratischer Nichtrest. Wir nennen dieses $\varepsilon$, von dem der biquadratische Charakter der Zahl $M$ bezüglich des Moduls $m$ abhängt, den Dezidenten von $M$ zum Modul $m$. Durch Induktion erhält man den folgenden schönen Satz. „Sind $M$ und $m$ ungerade Primzahlen der Form $1+(2+2i)\mu$, wobei $\mu$ eine ganze Zahl ist, so ist die Differenz der beiden Dezidenten, von $M$ zum Modul $m$ und von $m$ zum Modul $M$, entweder $=0$ oder $=2$; das Erste gilt, wenn wenigstens eine der Zahlen $m,M$ die Form $1+4N$ hat, das Andere, wenn beide die Form $1+2i+4N$ haben.“ Dieser Reziprozitätssatz ist dem für quadratische Reste bei ausschließlich reellen Zahlen analog.

\Article{4}
Man multipliziere alle Zahlen $f$ mit $M$ und bestimme ihre kleinsten Reste zum Modul $m$. Befinden sich darunter
\[
\begin{array}{cll}
\alpha&\text{zu }f&\text{gehörig}\\
\beta&f'&\\
\gamma&f''&\\
\delta&f'''&
\end{array}
\]
so ist $\varepsilon\equiv\beta+2\gamma+3\delta\ (\mathrm{mod.}\ 4)$.

\SourceLeaf{343}
Beweis. Die Gesamtheit derjenigen Zahlen aus $f$, deren Produkte mit $M$ Reste aus $f$ ergeben, sei $g$; die Gesamtheit derjenigen, deren Produkte Reste aus $f'$ ergeben, sei $g'$, und ebenso $g''\pdot g'''$. Dann sind die kleinsten Reste von
\begingroup\postdisplaypenalty=10000
\[
-ig'M,\quad -g''M,\quad ig'''M
\]
sämtlich in $f$ enthalten und sowohl voneinander als auch von den kleinsten Resten der Produkte $gM$ verschieden. Folglich ist das Produkt aller
\[
gM,\quad -ig'M,\quad -g''M,\quad +ig'''M
\]
zum Produkt aller $f$ kongruent, also auch zum Produkt aller $g,g',g'',g'''$. Jenes Produkt ist aber gleich dem Produkt aller $g,g',g'',g'''$ multipliziert mit
\[
M^{\alpha}\pdot(-iM)^{\beta}\pdot(-M)^{\gamma}\pdot(iM)^{\delta}
\]
also ist dieses letzte Produkt $\equiv1$
\[
\begin{aligned}
\text{folglich}\qquad&M^{\alpha+\beta+\gamma+\delta}(-i)^{\beta}(-1)^{\gamma}i^{\delta}\equiv1\\
\text{oder}\qquad&M^{\alpha+\beta+\gamma+\delta}\equiv i^{\beta}(-1)^{\gamma}(-i)^{\delta}\equiv i^{\beta+2\gamma+3\delta}
\end{aligned}
\]
woraus der Satz unmittelbar folgt.\endgroup

\Needspace{19\baselineskip}\Article{5}
Ob der kleinste Rest einer Zahl $N$ zum Modul $m$ zur Klasse $f,f',f''$ oder $f'''$ gehört, lässt sich leicht entscheiden. Ist nämlich $\omega$ die in $\frac{N}{m}$ enthaltene ganze Zahl, so ist jener Rest $=N-\omega m$ und gehört somit zu $f,f',f''$, je nachdem nach dem Ansatz
\[
\frac{N}{m}-\omega=x+iy
\]
die folgenden Bedingungen gelten:
\[
\begin{array}{ll}
x<\tfrac12,&y<\tfrac12\\
x>\tfrac12,&y<\tfrac12\\
x>\tfrac12,&y>\tfrac12\\
x<\tfrac12,&y>\tfrac12
\end{array}
\]
In diesen $4$ Fällen ist die in $\frac{2N}{m}$ enthaltene ganze Zahl der Reihe nach die folgende:
\SourceLeaf{344}
\[
\begin{array}{l}
2\omega\\
2\omega+1\\
2\omega+1+i\\
2\omega+i
\end{array}
\]
Daraus ist ersichtlich, dass der kleinste Rest von $N$ zum Modul $m$ zu $f,f',f'',f'''$ gehört, je nachdem die folgenden Bedingungen gelten, wenn man die in $\frac{2N}{m}$ enthaltene ganze Zahl $=\xi+\eta i$ setzt:
\[
\begin{array}{ll}
\xi\ \text{gerade}&\eta\ \text{gerade}\\
\xi\ \text{ungerade}&\eta\ \text{gerade}\\
\xi\ \text{ungerade}&\eta\ \text{ungerade}\\
\xi\ \text{gerade}&\eta\ \text{ungerade}
\end{array}
\]

\Needspace{27\baselineskip}\Article{6}
Danach erhält man den Dezidenten von $M$ zum Modul $m$ wie folgt. Man bestimme die ganzen Zahlen, die in den einzelnen $\frac{2fM}{m}$ enthalten sind. Bezeichnet man diese allgemein mit $x+yi$, so lasse man diejenigen ganz außer Acht, bei denen $x$ und $y$ beide gerade sind; man rechne für jede derjenigen, bei denen $x$ ungerade und $y$ gerade ist, eins, nehme für jede derjenigen, bei denen $x$ und $y$ beide ungerade sind, zwei und für jede derjenigen, bei denen $x$ gerade und $y$ ungerade ist, drei. Von der Summe aller dieser Zahlen nehme man den kleinsten Rest modulo $4$; dies ist der gesuchte Dezident. Wir drücken dies so aus:
\[
\mathrm{Dec.}\frac{M}{m}=\Sigma n
\]
wobei $[\frac{2fM}{m}]=x+yi$ und $n=0$ zu setzen ist, wenn $x$ gerade und $y$ gerade ist
\[
\begin{array}{cll}
1&x\ \text{ungerade}&y\ \text{gerade}\\
2&x\ \text{ungerade}&y\ \text{ungerade}\\
3&x\ \text{gerade}&y\ \text{ungerade}
\end{array}
\]
Der Kürze halber wollen wir $n$ mit der Charakteristik $\Theta$ bezeichnen, \mbox{$n=\Theta\frac{2fM}{m}$}\begingroup\renewcommand{\thefootnote}{*)}\footnote{Um zu entscheiden, zu welcher Klasse $M$ bezüglich $m$ gehört, wählt man diejenigen Werte von $k$ (unter den Zahlen $1,2,3\ldots p-1$) aus, für die $[\frac{2km'M}{p}]$ gerade wird, und addiert $-\Sigma[\frac{2km'}{p}]^2$. Nimmt man $k$ nur bis $\frac12p$, so ist zu summieren
\[
-\Sigma\{[\frac{2km'M}{p}]^2+[\frac{2km'}{p}]^2\}
\]
für diejenigen Werte von $[\frac{2km'M}{p}]$, die durch $1+i$ teilbar sind.}\endgroup.
